\documentclass[11pt,a4paper]{article}
\usepackage{iftex}
\ifPDFTeX
\usepackage[utf8]{inputenc}
\usepackage[T1]{fontenc}
\usepackage{lmodern}
\else
\usepackage{fontspec}
\setmainfont{lmroman10-regular.otf}[BoldFont=lmroman10-bold.otf,ItalicFont=lmroman10-italic.otf,BoldItalicFont=lmroman10-bolditalic.otf]
\setmonofont{lmmono10-regular.otf}
\fi
\usepackage[ngerman]{babel}
\usepackage[a4paper,margin=20mm]{geometry}
\usepackage{microtype}
\usepackage{amsmath,amssymb}
\usepackage{booktabs,longtable,tabularx,array}
\usepackage{enumitem}
\usepackage{xcolor}
\usepackage{xurl}
\usepackage{hyperref}

\definecolor{navy}{HTML}{12334A}
\definecolor{teal}{HTML}{146B75}
\definecolor{signal}{HTML}{A8391A}
\definecolor{softgrey}{HTML}{F1F3F5}
\hypersetup{colorlinks=true,linkcolor=navy,citecolor=teal,urlcolor=teal,
  pdftitle={Datenkatalog des S\&P-500-Nachhaltigkeitsrankings},
  pdfauthor={ETHack 2026}}
\urlstyle{same}

\newcolumntype{P}[1]{>{\raggedright\arraybackslash}p{#1}}
\newcolumntype{R}[1]{>{\raggedleft\arraybackslash}p{#1}}
\setlength{\parindent}{0pt}
\setlength{\parskip}{5pt}
\setlength{\tabcolsep}{4pt}
\setlength{\emergencystretch}{2em}
\setlist{itemsep=3pt,topsep=4pt,leftmargin=*}
\renewcommand{\arraystretch}{1.15}
\setcounter{tocdepth}{2}

\newcommand{\befund}[1]{%
  \par\vspace{2pt}\noindent\colorbox{softgrey}{%
  \begin{minipage}{\dimexpr\linewidth-2\fboxsep}\small\vspace{3pt}#1\vspace{3pt}\end{minipage}}\par\vspace{4pt}}

\title{\textcolor{navy}{Datenkatalog}\\[4mm]
\large Alles, was für das S\&P-500-Nachhaltigkeitsranking gesammelt wurde --\\
einmal nach Quelle, einmal nach Datentyp}
\author{ETHack 2026}
\date{Stand aller Abrufe: 12.~September 2026}

\begin{document}
\maketitle

\befund{\textbf{Wozu dieser Katalog.} Der Bestand ist an einem Tag aus zwölf Quellen zusammengetragen worden. Damit er später wiederverwendbar ist, braucht es zwei Dinge: eine Liste, woher jeder Wert stammt und was sein Abruf kostet, und eine Liste, was die Werte eigentlich messen. Beide Tabellen sind \emph{erzeugt}, nicht geschrieben -- \texttt{scripts/06\_katalog.py} liest sie aus den Daten. Kommt eine Quelle dazu, stimmt der Katalog beim nächsten Lauf von selbst.}

\section*{Der Bestand in Zahlen}
\addcontentsline{toc}{section}{Der Bestand in Zahlen}

\begin{center}\small
\begin{tabular}{lr}
\toprule
Quellen & 12 \\
Rohtabellen im Zwischenspeicher & 21 \\
Rohzeilen insgesamt & 9\,750\,265 \\
Umfang der Rohdaten & 253{,}7 MB \\
\midrule
Werte im gesammelten Datensatz & 7\,403 \\
davon Firmen & 436 von 503 \\
davon Kennzahlen & 27 \\
abgedeckte Geschäftsjahre & 2018--2026 \\
\midrule
Werte ohne Quellenangabe & \textbf{0} \\
Quellen mit Schlüsselpflicht & 2 \\
\bottomrule
\end{tabular}
\end{center}

\tableofcontents
\newpage

\section{Aufbau des gesammelten Datensatzes}

Der Datensatz liegt im Langformat vor: \textbf{eine Zeile je Firma, Jahr und Kennzahl}. Das ist absichtlich so und nicht als breite Tabelle, weil sich nur so die Herkunft an jedem einzelnen \emph{Wert} führen lässt statt nur an der Spalte.

Der Unterschied ist nicht kosmetisch. Die CO\textsubscript{2}-Zahl einer Firma kann am Schornstein gemessen sein, die der nächsten aus Brennstoffmengen gerechnet und die der dritten aus einem Geschäftsbericht abgeschrieben. Eine breite Tabelle würde diese drei Fälle in einer Spalte \texttt{co2} zusammenwerfen und die Unterscheidung verlieren -- genau die Information, an der ein Ranking steht oder fällt.

\begin{center}\small
\begin{tabular}{lP{10.6cm}}
\toprule
\textbf{Spalte} & \textbf{Inhalt} \\
\midrule
\texttt{ticker}, \texttt{company} & Firma \\
\texttt{gics\_sector}, \texttt{gics\_sub\_industry} & Branchenzuordnung \\
\texttt{year} & Geschäftsjahr des Wertes \\
\texttt{metric}, \texttt{metric\_label} & Kennzahl, technisch und lesbar \\
\texttt{value}, \texttt{unit} & Wert und Einheit \\
\texttt{direction} & $-1$ kleiner ist besser, $+1$ größer ist besser, $0$ neutral \\
\texttt{axis} & A Kernnote, B Glaubwürdigkeit, S sozial, G Governance, sonst Bezugsgröße \\
\texttt{source\_id}, \texttt{source\_name} & Quelle \\
\texttt{source\_url} & Abrufadresse \\
\texttt{source\_access}, \texttt{source\_license} & Zugangsbedingung und Lizenz \\
\texttt{retrieved\_at} & Abrufdatum \\
\bottomrule
\end{tabular}
\end{center}

\textbf{Dateien.} \texttt{data/out/dataset\_long.parquet} und \texttt{.csv} enthalten alle Werte. \texttt{sources.json} und \texttt{metrics.json} sind die beiden Register. \texttt{companies.json} ist eine je Firma verdichtete Sicht, die die Web-Oberfläche direkt liest.

\newpage
\section{Sicht 1: nach Quelle}

Je Quelle steht hier: ob ein Schlüssel nötig ist, unter welcher Lizenz sie steht, wie viele Rohzeilen im Zwischenspeicher liegen, für wie viele Firmen sie am Ende einen Wert liefert -- und wo ihre Grenze liegt.

\begin{center}\footnotesize
\begin{longtable}{P{5.4cm} c P{2.9cm} r r r}
\toprule
\textbf{Quelle} & \textbf{Key} & \textbf{Lizenz} & \textbf{Rohzeilen} & \textbf{Firmen} & \textbf{Werte} \\
\midrule
\endfirsthead
\toprule
\textbf{Quelle} & \textbf{Key} & \textbf{Lizenz} & \textbf{Rohzeilen} & \textbf{Firmen} & \textbf{Werte} \\
\midrule
\endhead
\textbf{EPA Greenhouse Gas Reporting Program} & nein & US-Regierungswerk, gemeinfrei & 130\,036 & 136 & 1\,017 \\
\multicolumn{6}{p{0.96\textwidth}}{\footnotesize\ttfamily https:/\allowbreak{}/\allowbreak{}data.\allowbreak{}epa.\allowbreak{}gov/\allowbreak{}efservice/\allowbreak{}PUB\_DIM\_FACILITY/\allowbreak{}}\\
\multicolumn{6}{p{0.96\textwidth}}{\footnotesize US-Anlagen ueber 25 000 t CO2e, Berichtsjahre 2018-2023 \quad\textit{Grenze:} Reihe endet 2023; Direktemittenten von Lieferanten getrennt (sector\_type)}\\
\addlinespace[3pt]
\textbf{EPA Clean Air Markets Program Data} & ja & US-Regierungswerk, gemeinfrei & 21\,912 & 33 & 499 \\
\multicolumn{6}{p{0.96\textwidth}}{\footnotesize\ttfamily https:/\allowbreak{}/\allowbreak{}api.\allowbreak{}epa.\allowbreak{}gov/\allowbreak{}easey/\allowbreak{}emissions-\allowbreak{}mgmt/\allowbreak{}emissions/\allowbreak{}apportioned/\allowbreak{}annual}\\
\multicolumn{6}{p{0.96\textwidth}}{\footnotesize Kraftwerksbloecke der Handelsprogramme, 2019 bis laufendes Jahr \quad\textit{Grenze:} nur Stromerzeugung; laufendes Jahr unvollstaendig}\\
\addlinespace[3pt]
\textbf{EPA eGRID} & nein & US-Regierungswerk, gemeinfrei & 12\,239 & 43 & 69 \\
\multicolumn{6}{p{0.96\textwidth}}{\footnotesize\ttfamily https:/\allowbreak{}/\allowbreak{}www.\allowbreak{}epa.\allowbreak{}gov/\allowbreak{}egrid/\allowbreak{}download-\allowbreak{}data}\\
\multicolumn{6}{p{0.96\textwidth}}{\footnotesize US-Kraftwerke, Berichtsjahr 2023 \quad\textit{Grenze:} nur aussagekraeftig fuer Firmen, deren Geschaeft Stromerzeugung ist}\\
\addlinespace[3pt]
\textbf{SEC EDGAR XBRL frames} & nein & oeffentlich & 53\,492 & 133 & 732 \\
\multicolumn{6}{p{0.96\textwidth}}{\footnotesize\ttfamily https:/\allowbreak{}/\allowbreak{}data.\allowbreak{}sec.\allowbreak{}gov/\allowbreak{}api/\allowbreak{}xbrl/\allowbreak{}frames/\allowbreak{}}\\
\multicolumn{6}{p{0.96\textwidth}}{\footnotesize Geschaeftsjahre 2016-2025 \quad\textit{Grenze:} fuenf verschiedene Umsatz-Tags; Firmen taggen uneinheitlich}\\
\addlinespace[3pt]
\textbf{Science Based Targets initiative} & nein & Nutzung mit Quellenangabe & 15\,605 & 239 & 1\,228 \\
\multicolumn{6}{p{0.96\textwidth}}{\footnotesize\ttfamily https:/\allowbreak{}/\allowbreak{}sciencebasedtargets.\allowbreak{}org/\allowbreak{}download/\allowbreak{}excel}\\
\multicolumn{6}{p{0.96\textwidth}}{\footnotesize 15 605 Organisationen weltweit, Stand laufend \quad\textit{Grenze:} Selbstselektion: Firmen melden sich freiwillig; Energiebranche fehlt fast ganz}\\
\addlinespace[3pt]
\textbf{EPA Toxics Release Inventory} & nein & US-Regierungswerk, gemeinfrei & 36\,108 & 194 & 582 \\
\multicolumn{6}{p{0.96\textwidth}}{\footnotesize\ttfamily https:/\allowbreak{}/\allowbreak{}data.\allowbreak{}epa.\allowbreak{}gov/\allowbreak{}efservice/\allowbreak{}downloads/\allowbreak{}tri/\allowbreak{}mv\_tri\_basic\_download/\allowbreak{}2023\_US/\allowbreak{}csv/\allowbreak{}}\\
\multicolumn{6}{p{0.96\textwidth}}{\footnotesize US-Anlagen, Berichtsjahr 2023 \quad\textit{Grenze:} Gesamtmenge korreliert mit CO2 (rho 0,53); eigenstaendig ist der Krebserreger-Anteil}\\
\addlinespace[3pt]
\textbf{OSHA Injury Tracking Application, Formular 300A} & nein & US-Regierungswerk, gemeinfrei & 1\,176\,137 & 313 & 939 \\
\multicolumn{6}{p{0.96\textwidth}}{\footnotesize\ttfamily https:/\allowbreak{}/\allowbreak{}www.\allowbreak{}osha.\allowbreak{}gov/\allowbreak{}itadata}\\
\multicolumn{6}{p{0.96\textwidth}}{\footnotesize US-Betriebsstaetten, 2023-2025 \quad\textit{Grenze:} DART misst auch Branchenstruktur; branchenrelativ vergleichen}\\
\addlinespace[3pt]
\textbf{EPA ECHO Enforcement and Compliance History} & nein & US-Regierungswerk, gemeinfrei & 3\,178\,967 & 134 & 402 \\
\multicolumn{6}{p{0.96\textwidth}}{\footnotesize\ttfamily https:/\allowbreak{}/\allowbreak{}echo.\allowbreak{}epa.\allowbreak{}gov/\allowbreak{}files/\allowbreak{}echodownloads/\allowbreak{}echo\_exporter.\allowbreak{}zip}\\
\multicolumn{6}{p{0.96\textwidth}}{\footnotesize US-Anlagen, Konformitaetshistorie drei Jahre \quad\textit{Grenze:} Zuordnung ueber FRS-Kennung, deshalb nur fuer Anlagen mit GHGRP-Eintrag}\\
\addlinespace[3pt]
\textbf{EIA Formular 923 (via PUDL und API)} & ja & US-Regierungswerk, gemeinfrei & 264\,315 & -- & -- \\
\multicolumn{6}{p{0.96\textwidth}}{\footnotesize\ttfamily https:/\allowbreak{}/\allowbreak{}www.\allowbreak{}eia.\allowbreak{}gov/\allowbreak{}electricity/\allowbreak{}data/\allowbreak{}eia923/\allowbreak{}}\\
\multicolumn{6}{p{0.96\textwidth}}{\footnotesize US-Kraftwerke ab 1 MW, 2021-2025 \quad\textit{Grenze:} unabhaengige Gegenprobe zu CEMS, nicht dieselbe Messkette}\\
\addlinespace[3pt]
\textbf{Kommerzielle ESG-Risk-Ratings (Sustainalytics via Yahoo, gespiegelt)} & nein & unklar -- nur als Vergleichsma & 245 & 199 & 199 \\
\multicolumn{6}{p{0.96\textwidth}}{\footnotesize\ttfamily https:/\allowbreak{}/\allowbreak{}raw.\allowbreak{}githubusercontent.\allowbreak{}com/\allowbreak{}sburstein/\allowbreak{}ESG-\allowbreak{}Stock-\allowbreak{}Data/\allowbreak{}main/\allowbreak{}sp\_esg\_stock\_data.\allowbreak{}csv}\\
\multicolumn{6}{p{0.96\textwidth}}{\footnotesize 245 Firmen, Stand 2021 \quad\textit{Grenze:} eingefroren 2021; fliesst in keine eigene Note ein}\\
\addlinespace[3pt]
\textbf{S\&P-500-Konstituenten} & nein & CC BY-SA & 503 & -- & -- \\
\multicolumn{6}{p{0.96\textwidth}}{\footnotesize\ttfamily https:/\allowbreak{}/\allowbreak{}en.\allowbreak{}wikipedia.\allowbreak{}org/\allowbreak{}wiki/\allowbreak{}List\_of\_S\%26P\_500\_companies}\\
\multicolumn{6}{p{0.96\textwidth}}{\footnotesize 503 Mitglieder, heutiger Stand \quad\textit{Grenze:} nur der heutige Stand -- fuer historische Jahre ueberlebensverzerrt}\\
\addlinespace[3pt]
\textbf{PUDL (Catalyst Cooperative) auf Zenodo} & nein & Public Domain & 1\,162\,085 & -- & -- \\
\multicolumn{6}{p{0.96\textwidth}}{\footnotesize\ttfamily https:/\allowbreak{}/\allowbreak{}s3.\allowbreak{}us-\allowbreak{}west-\allowbreak{}2.\allowbreak{}amazonaws.\allowbreak{}com/\allowbreak{}pudl.\allowbreak{}catalyst.\allowbreak{}coop/\allowbreak{}stable/\allowbreak{}}\\
\multicolumn{6}{p{0.96\textwidth}}{\footnotesize aufbereitete EIA-, EPA- und SEC-Daten, 1995 bis laufendes Jahr \quad\textit{Grenze:} Zwischenschicht: Aufbereitungsfehler sind moeglich, dafuer entfaellt eigene Bereinigung}\\
\addlinespace[3pt]
\textbf{SEC 10-K Anlage 21 (Konzernstruktur, via PUDL)} & nein & Public Domain & 3\,800\,706 & -- & -- \\
\multicolumn{6}{p{0.96\textwidth}}{\footnotesize\ttfamily https:/\allowbreak{}/\allowbreak{}s3.\allowbreak{}us-\allowbreak{}west-\allowbreak{}2.\allowbreak{}amazonaws.\allowbreak{}com/\allowbreak{}pudl.\allowbreak{}catalyst.\allowbreak{}coop/\allowbreak{}stable/\allowbreak{}core\_sec10k\_\_quarterly\_exhibit\_21\_company\_ownership.\allowbreak{}parquet}\\
\multicolumn{6}{p{0.96\textwidth}}{\footnotesize 3,8 Mio. Tochtereintraege, 478 von 500 Indexmuettern \quad\textit{Grenze:} Namen in Kleinschreibung und teils abgeschnitten; Beteiligungsquote nur bei 99 117 Eintraegen}\\
\addlinespace[3pt]
\textbf{SEC 10-K-Finanzkennzahlen (companyfacts, aufbereitet von Janis)} & nein & oeffentlich & 503 & 495 & 2\,269 \\
\multicolumn{6}{p{0.96\textwidth}}{\footnotesize\ttfamily https:/\allowbreak{}/\allowbreak{}data.\allowbreak{}sec.\allowbreak{}gov/\allowbreak{}api/\allowbreak{}xbrl/\allowbreak{}companyfacts/\allowbreak{}}\\
\multicolumn{6}{p{0.96\textwidth}}{\footnotesize Geschaeftsjahr 2024, 499 Indexfirmen \quad\textit{Grenze:} Schulden uneinheitlich getaggt (debt\_tag); F\&E nur bei 135 Firmen berichtet}\\
\addlinespace[3pt]
\textbf{DOL Wage and Hour Division, Verfahren (aufbereitet von Janis)} & nein & US-Regierungswerk, gemeinfrei & 29\,374 & 139 & 695 \\
\multicolumn{6}{p{0.96\textwidth}}{\footnotesize\ttfamily https:/\allowbreak{}/\allowbreak{}data.\allowbreak{}dol.\allowbreak{}gov/\allowbreak{}data-\allowbreak{}catalog/\allowbreak{}WHD/\allowbreak{}enforcement/\allowbreak{}WHD\_enforcement.\allowbreak{}zip}\\
\multicolumn{6}{p{0.96\textwidth}}{\footnotesize Feststellungsende 2022-2024 \quad\textit{Grenze:} kein CIK, Zuordnung ueber Namen; Firmen ohne Treffer fehlen statt als Null zu erscheinen. Gegenprobe mit eigenem exakten Abgleich: 127 statt 213 Firmen, Rangkorrelation der Fallzahlen 0,62 -- die Team-Zuordnung zaehlt vermutlich Franchise-Betriebe unter dem Markennamen mit (McDonald's 142 gegen 64 Faelle)}\\
\addlinespace[3pt]
\textbf{SEC Form SD, Konfliktmineralien (aufbereitet von Janis)} & nein & oeffentlich & -- & 500 & 500 \\
\multicolumn{6}{p{0.96\textwidth}}{\footnotesize\ttfamily https:/\allowbreak{}/\allowbreak{}www.\allowbreak{}sec.\allowbreak{}gov/\allowbreak{}cgi-\allowbreak{}bin/\allowbreak{}browse-\allowbreak{}edgar?\allowbreak{}action=getcompany\&\allowbreak{}type=SD}\\
\multicolumn{6}{p{0.96\textwidth}}{\footnotesize Meldungen 2022-2025 \quad\textit{Grenze:} zeigt nur, dass gemeldet wird -- nicht, was die Meldung ueber Schmelzen und Herkunft sagt}\\
\addlinespace[3pt]
\textbf{World Benchmarking Alliance, Unternehmensprofile 2026 (erhoben von Janis)} & nein & CC BY 4.0 -- Namensnennung: Wo & 218 & 217 & 939 \\
\multicolumn{6}{p{0.96\textwidth}}{\footnotesize\ttfamily https:/\allowbreak{}/\allowbreak{}www.\allowbreak{}worldbenchmarkingalliance.\allowbreak{}org/\allowbreak{}company-\allowbreak{}scoreboard}\\
\multicolumn{6}{p{0.96\textwidth}}{\footnotesize 217 Indexfirmen, Bewertungsrunde 2026 \quad\textit{Grenze:} haengt an Offenlegung; Benchmarks nicht unabhaengig; CTT fast ohne Streuung (216 von 217 = 0)}\\
\addlinespace[3pt]
\textbf{Wikidata, offizielle Websites (fuer Firmenlogos)} & nein & CC0 & -- & -- & -- \\
\multicolumn{6}{p{0.96\textwidth}}{\footnotesize\ttfamily https:/\allowbreak{}/\allowbreak{}query.\allowbreak{}wikidata.\allowbreak{}org/\allowbreak{}sparql}\\
\multicolumn{6}{p{0.96\textwidth}}{\footnotesize Website ueber CIK (P5531 -> P856), Rest aus SEC Submissions \quad\textit{Grenze:} nur fuer die Anzeige; Logos sind Favicons und Marken der jeweiligen Firmen}\\
\addlinespace[3pt]
\textbf{Eigene Berechnung dieser Pipeline} & nein & - & -- & 133 & 1\,724 \\
\multicolumn{6}{p{0.96\textwidth}}{\footnotesize\ttfamily scripts/\allowbreak{}02\_analyse.\allowbreak{}py}\\
\multicolumn{6}{p{0.96\textwidth}}{\footnotesize abgeleitet aus den oben genannten Quellen \quad\textit{Grenze:} Herkunft der Eingangsgroessen steht jeweils beim Basiswert}\\
\bottomrule%
\end{longtable}
\end{center}

\befund{\textbf{Nur zwei Quellen verlangen eine Anmeldung}, beide kostenlos und sofort. Der EPA-Schlüssel für die Clean Air Markets ist der wichtige: Ohne ihn greift \texttt{DEMO\_KEY} mit zehn Anfragen je Stunde, mit ihm sind es tausend. Er ist auch der Schlüssel, der die zeitliche Lücke nach 2023 schließt. Der EIA-Schlüssel dient nur der unabhängigen Gegenprobe. Beide liegen in \texttt{.env}, die in \texttt{.gitignore} steht; eine Vorlage ohne Werte liegt als \texttt{.env.example} bei.}

\subsection{Rohdaten im Zwischenspeicher}

Alle Abrufe werden als Parquet zwischengespeichert. Das macht die Auswertung reproduzierbar und schont die Server -- zwei der Quellen brauchen mehrere Minuten je Abruf.

\begin{center}\footnotesize
\begin{longtable}{P{7.2cm} r r r c}
\toprule
\textbf{Tabelle} & \textbf{Zeilen} & \textbf{Spalten} & \textbf{MB} & \textbf{Jahre} \\
\midrule
\endfirsthead
\toprule
\textbf{Tabelle} & \textbf{Zeilen} & \textbf{Spalten} & \textbf{MB} & \textbf{Jahre} \\
\midrule
\endhead
\texttt{sec10k\_ex21} & 3\,800\,706 & 5 & 76.3 & -- \\
\texttt{epa\_echo} & 3\,178\,967 & 10 & 92.3 & -- \\
\texttt{osha\_ita} & 1\,176\,137 & 17 & 53.1 & 2023--2025 \\
\texttt{eia923\_gen\_fuel} & 727\,126 & 9 & 12.4 & 2001--2025 \\
\texttt{eia\_generation} & 181\,340 & 4 & 2.5 & 2001--2025 \\
\texttt{out\_eia\_\_yearly\_utilities} & 165\,728 & 27 & 2.3 & 2001--2026 \\
\texttt{out\_epacems\_\_yearly\_operational\_characteristics} & 118\,117 & 13 & 3.8 & 2000--2025 \\
\texttt{out\_eia860\_\_yearly\_ownership} & 106\,281 & 16 & 0.7 & 2001--2025 \\
\texttt{eia\_api\_facility\_fuel} & 72\,259 & 6 & 1.3 & 2019--2025 \\
\texttt{epa\_facility} & 65\,155 & 10 & 1.2 & 2018--2023 \\
\texttt{epa\_emission} & 64\,881 & 5 & 0.6 & 2018--2023 \\
\texttt{sec\_revenue} & 53\,492 & 5 & 0.8 & 2016--2025 \\
\texttt{core\_epa\_\_assn\_eia\_epacamd} & 44\,833 & 7 & 0.2 & 2018--2024 \\
\texttt{epa\_tri} & 36\,108 & 122 & 4.8 & 2023--2023 \\
\texttt{sbti\_targets} & 15\,605 & 13 & 0.5 & -- \\
\texttt{egrid\_plant} & 12\,239 & 11 & 0.6 & 2023--2023 \\
\texttt{campd\_facility} & 11\,051 & 8 & 0.1 & 2019--2026 \\
\texttt{campd\_emission} & 10\,861 & 7 & 0.2 & 2019--2026 \\
\texttt{epa\_eia\_crosswalk} & 10\,716 & 3 & 0.0 & 2018--2024 \\
\texttt{sp500\_master} & 503 & 5 & 0.0 & -- \\
\texttt{esg\_snapshot} & 245 & 11 & 0.0 & -- \\
\bottomrule%
\end{longtable}
\end{center}

\newpage
\section{Sicht 2: nach Datentyp}

Dieselben Werte, gruppiert nach dem, was sie messen. Der Pfeil gibt die Richtung an: $\downarrow$ heißt, kleiner ist besser.

\begin{center}\footnotesize
\begin{longtable}{P{4.1cm} P{3.5cm} P{2.9cm} c r c P{2.7cm}}
\toprule
\textbf{Kennzahl} & \textbf{technisch} & \textbf{Einheit} & \textbf{Richt.} & \textbf{Firmen} & \textbf{Jahre} & \textbf{Quelle} \\
\midrule
\endfirsthead
\toprule
\textbf{Kennzahl} & \textbf{technisch} & \textbf{Einheit} & \textbf{Richt.} & \textbf{Firmen} & \textbf{Jahre} & \textbf{Quelle} \\
\midrule
\endhead
\textbf{Emissionsmengen} & & & & & & & \\
\multicolumn{8}{p{0.96\textwidth}}{\footnotesize Absolute Treibhausgasmengen, wie sie gemeldet oder gemessen wurden.}\\
\addlinespace[2pt]
Scope-1-Emissionen & \texttt{scope1\_t} & t CO2e & $\downarrow$ & aggregiert & 136 & 2018--2023 & EPA Greenhouse Gas Reporti \\
CO2 gemessen (CEMS) & \texttt{campd\_co2\_t} & t CO2e & $\downarrow$ & aggregiert & 32 & 2019--2026 & EPA Clean Air Markets Prog \\
\midrule
\midrule
\textbf{Kraftwerksbilanz} & & & & & & & \\
\multicolumn{8}{p{0.96\textwidth}}{\footnotesize Menge und Erzeugung je Firma, beide von eGRID gemeldet. Eine Rate je MWh steht hier bewusst nicht: Die waere unsere Rechnung, nicht die der Quelle.}\\
\addlinespace[2pt]
CO2 der Kraftwerke & \texttt{egrid\_co2\_t} & t CO2e & $\downarrow$ & aggregiert & 43 & 2023--2023 & EPA eGRID \\
Nettoerzeugung & \texttt{egrid\_mwh} & MWh & -- & aggregiert & 43 & 2023--2023 & EPA eGRID \\
eGRID-Kraftwerke & \texttt{egrid\_plants} & Anzahl & -- & aggregiert & 43 & 2023--2023 & EPA eGRID \\
\midrule
\midrule
\textbf{Weitere Umweltwirkung} & & & & & & & \\
\multicolumn{8}{p{0.96\textwidth}}{\footnotesize Was neben CO2 in Luft, Wasser und Boden gelangt.}\\
\addlinespace[2pt]
Giftstofffreisetzung & \texttt{tri\_releases\_lbs} & lbs & $\downarrow$ & aggregiert & 194 & 2023--2023 & EPA Toxics Release Invento \\
davon krebserregend & \texttt{tri\_carcinogen\_lbs} & lbs & $\downarrow$ & aggregiert & 194 & 2023--2023 & EPA Toxics Release Invento \\
\midrule
\midrule
\textbf{Arbeitssicherheit} & & & & & & & \\
\multicolumn{8}{p{0.96\textwidth}}{\footnotesize Die soziale Dimension als gemeldete Kennzahl statt als Selbstbeschreibung.}\\
\addlinespace[2pt]
Faelle mit Ausfalltagen & \texttt{osha\_dafw\_cases} & Anzahl & $\downarrow$ & aggregiert & 314 & 2025--2025 & OSHA Injury Tracking Appli \\
Faelle mit Einschraenkung & \texttt{osha\_djtr\_cases} & Anzahl & $\downarrow$ & aggregiert & 314 & 2025--2025 & OSHA Injury Tracking Appli \\
gemeldete Arbeitsstunden & \texttt{osha\_hours} & Stunden & -- & aggregiert & 314 & 2025--2025 & OSHA Injury Tracking Appli \\
Todesfaelle & \texttt{osha\_deaths} & Anzahl & $\downarrow$ & aggregiert & 314 & 2025--2025 & OSHA Injury Tracking Appli \\
gemeldete Betriebsstaetten & \texttt{osha\_sites} & Anzahl & -- & aggregiert & 314 & 2025--2025 & OSHA Injury Tracking Appli \\
\midrule
\midrule
\textbf{Regeltreue} & & & & & & & \\
\multicolumn{8}{p{0.96\textwidth}}{\footnotesize Dokumentiertes Verhalten gegenueber Umweltauflagen. Die einzige Achse, die nachweislich unabhaengig von der CO2-Intensitaet ist.}\\
\addlinespace[2pt]
Umweltstrafen & \texttt{echo\_penalties\_usd} & USD & $\downarrow$ & aggregiert & 134 & 2025--2025 & EPA ECHO Enforcement and C \\
Verstossquartale & \texttt{echo\_nc\_quarters} & Quartale & $\downarrow$ & aggregiert & 134 & 2025--2025 & EPA ECHO Enforcement and C \\
ECHO-Anlagen & \texttt{echo\_facilities} & Anzahl & -- & aggregiert & 134 & 2025--2025 & EPA ECHO Enforcement and C \\
Anlagen mit schwerem Verstoss & \texttt{echo\_significant} & Anzahl & $\downarrow$ & aggregiert & 134 & 2025--2025 & EPA ECHO Enforcement and C \\
\midrule
\midrule
\textbf{Ziele und Glaubwuerdigkeit} & & & & & & & \\
\multicolumn{8}{p{0.96\textwidth}}{\footnotesize Achse B des Modells. Wird getrennt ausgewiesen und nie in die Kernnote eingerechnet.}\\
\addlinespace[2pt]
SBTi-geprueftes Ziel & \texttt{sbti\_validated} & ja/nein & $\uparrow$ & gemeldet & 239 & 2026--2026 & Science Based Targets init \\
Zwischenzieljahr & \texttt{sbti\_near\_term\_year} & Jahr & -- & gemeldet & 207 & 2026--2026 & Science Based Targets init \\
Netto-null-Jahr & \texttt{sbti\_net\_zero\_year} & Jahr & -- & gemeldet & 65 & 2026--2026 & Science Based Targets init \\
Nahziel-Zusage zurueckgezogen & \texttt{sbti\_commitment\_removed} & ja/nein & $\downarrow$ & gemeldet & 239 & 2026--2026 & Science Based Targets init \\
Netto-null-Zusage zurueckgezogen & \texttt{sbti\_net\_zero\_removed} & ja/nein & $\downarrow$ & gemeldet & 239 & 2026--2026 & Science Based Targets init \\
Zwischenziel abgelaufen & \texttt{sbti\_near\_term\_expired} & ja/nein & $\downarrow$ & gemeldet & 239 & 2026--2026 & Science Based Targets init \\
\midrule
\midrule
\textbf{Arbeitsrecht} & & & & & & & \\
\multicolumn{8}{p{0.96\textwidth}}{\footnotesize Behoerdlich festgestellte Lohnverstoesse. Firmen ohne zuordenbaren Fall fehlen, statt als Null zu erscheinen.}\\
\addlinespace[2pt]
Lohnverfahren 2022-2024 & \texttt{whd\_cases} & Anzahl & $\downarrow$ & aggregiert & 139 & 2024--2024 & DOL Wage and Hour Division \\
festgestellte Verstoesse & \texttt{whd\_violations} & Anzahl & $\downarrow$ & aggregiert & 139 & 2024--2024 & DOL Wage and Hour Division \\
nachgezahlte Loehne & \texttt{whd\_backwages\_usd} & USD & $\downarrow$ & aggregiert & 139 & 2024--2024 & DOL Wage and Hour Division \\
betroffene Beschaeftigte & \texttt{whd\_employees} & Anzahl & $\downarrow$ & aggregiert & 139 & 2024--2024 & DOL Wage and Hour Division \\
Geldbussen Lohnrecht & \texttt{whd\_penalties\_usd} & USD & $\downarrow$ & aggregiert & 139 & 2024--2024 & DOL Wage and Hour Division \\
\midrule
\midrule
\textbf{Externe Bewertungen, offen lizenziert} & & & & & & & \\
\multicolumn{8}{p{0.96\textwidth}}{\footnotesize WBA-Benchmarks unter CC BY 4.0. Anders als der kommerzielle Snapshot ist die Methodik einsehbar -- bewertet wird trotzdem Offenlegung, nicht Wirkung.}\\
\addlinespace[2pt]
WBA Qualitaet des Transitionsplans & \texttt{wba\_tpq} & 0-5 & $\uparrow$ & gemeldet & 217 & 2026--2026 & World Benchmarking Allianc \\
WBA Beitrag zur Transition & \texttt{wba\_ctt} & 0-2 & $\uparrow$ & gemeldet & 217 & 2026--2026 & World Benchmarking Allianc \\
WBA Social Benchmark & \texttt{wba\_social} & 0-100 & $\uparrow$ & gemeldet & 217 & 2026--2026 & World Benchmarking Allianc \\
WBA Nature Benchmark & \texttt{wba\_nature} & 0-100 & $\uparrow$ & gemeldet & 71 & 2026--2026 & World Benchmarking Allianc \\
WBA Just Transition & \texttt{wba\_just\_transition} & 0-100 & $\uparrow$ & gemeldet & 217 & 2026--2026 & World Benchmarking Allianc \\
\midrule
\midrule
\textbf{Finanzkennzahlen} & & & & & & & \\
\multicolumn{8}{p{0.96\textwidth}}{\footnotesize Aus den 10-K-Berichten, als Bezugsgroessen fuer Intensitaeten und fuer die Portfoliofrage -- keine Nachhaltigkeitsbewertung.}\\
\addlinespace[2pt]
Nettogewinn & \texttt{net\_income\_usd} & USD & -- & gemeldet & 444 & 2024--2024 & SEC 10-K-Finanzkennzahlen \\
Bilanzsumme & \texttt{total\_assets\_usd} & USD & -- & gemeldet & 495 & 2024--2024 & SEC 10-K-Finanzkennzahlen \\
Finanzschulden & \texttt{total\_debt\_usd} & USD & -- & gemeldet & 391 & 2024--2024 & SEC 10-K-Finanzkennzahlen \\
operativer Cashflow & \texttt{operating\_cf\_usd} & USD & -- & gemeldet & 479 & 2024--2024 & SEC 10-K-Finanzkennzahlen \\
Sachinvestitionen & \texttt{capex\_usd} & USD & -- & gemeldet & 327 & 2024--2024 & SEC 10-K-Finanzkennzahlen \\
Forschung und Entwicklung & \texttt{rnd\_usd} & USD & -- & gemeldet & 133 & 2024--2024 & SEC 10-K-Finanzkennzahlen \\
\midrule
\midrule
\textbf{Lieferkette} & & & & & & & \\
\multicolumn{8}{p{0.96\textwidth}}{\footnotesize Ob eine Firma Konfliktmineralien meldet. Die Meldepflicht allein ist noch kein Befund ueber die Herkunft.}\\
\addlinespace[2pt]
meldet Konfliktmineralien (Form SD) & \texttt{sd\_conflict\_minerals\_filer} & ja/nein & -- & gemeldet & 500 & 2025--2025 & SEC Form SD, Konfliktminer \\
\midrule
\midrule
\textbf{Vergleichsmassstab} & & & & & & & \\
\multicolumn{8}{p{0.96\textwidth}}{\footnotesize Fremdbewertung, ausschliesslich zum Gegenhalten. Fliesst in keine eigene Note ein.}\\
\addlinespace[2pt]
kommerzielles ESG-Risiko & \texttt{esg\_risk\_total} & Punkte & $\downarrow$ & gemeldet & 199 & 2021--2021 & Kommerzielle ESG-Risk-Rati \\
\midrule
\midrule
\textbf{Bezugs- und Guetegroessen} & & & & & & & \\
\multicolumn{8}{p{0.96\textwidth}}{\footnotesize Nenner, Zaehlwerte und Konfidenzangaben, die keine Bewertung sind.}\\
\addlinespace[2pt]
Umsatz & \texttt{revenue\_musd} & Mio. USD & -- & gemeldet & 133 & 2018--2023 & SEC EDGAR XBRL frames \\
zugeordnete Anlagen & \texttt{n\_facilities} & Anzahl & -- & aggregiert & 136 & 2018--2023 & EPA Greenhouse Gas Reporti \\
Kraftwerke mit Messung & \texttt{campd\_plants} & Anzahl & -- & aggregiert & 33 & 2019--2026 & EPA Clean Air Markets Prog \\
TRI-Anlagen & \texttt{tri\_facilities} & Anzahl & -- & aggregiert & 194 & 2023--2023 & EPA Toxics Release Invento \\
\midrule
\bottomrule%
\end{longtable}
\end{center}

\section{Was beim Weiterverwenden zu beachten ist}

\subsection{Die Kennzahlen sind nicht gleichwertig}

Drei Arten von Werten stehen nebeneinander und dürfen nicht vermischt werden:

\begin{itemize}
\item \textbf{Gemessen.} CAMPD misst am Schornstein, ECHO und OSHA sind behördlich festgestellt beziehungsweise gesetzlich vorgeschrieben.
\item \textbf{Gemeldet.} GHGRP, TRI und SBTi beruhen auf Selbstauskunft -- teils geprüft, teils nicht.
\item \textbf{Berechnet.} Alles mit der Quelle \glqq Eigene Berechnung\grqq{} ist abgeleitet. Die Herkunft der Eingangsgrößen steht jeweils beim Basiswert.
\end{itemize}

\subsection{Zwei Quellen mit Selbstselektion}

SBTi und der kommerzielle ESG-Snapshot enthalten nur Firmen, die sich gemeldet haben beziehungsweise die der Anbieter bewertet. Dass die Energiebranche bei SBTi zu null Prozent vertreten ist, ist selbst eine Information -- aber keine, die man als Bewertung verrechnen darf.

\subsection{Eine Lizenz, die Vorsicht verlangt}

Der ESG-Snapshot ist über ein öffentliches Spiegel-Repositorium bezogen; die ursprüngliche Lizenzlage ist unklar. Er wird deshalb ausschließlich als Vergleichsmaßstab verwendet und fließt in keine eigene Note ein. Für eine Weitergabe des Datensatzes sollte er ausgenommen werden.

\subsection{Was der Katalog nicht enthält}

67 der 503 Indexmitglieder haben zu keiner Kennzahl einen Wert, überwiegend Finanzdienstleister. Sie stehen in \texttt{companies.json} unter \texttt{withoutData} -- ausdrücklich als \glqq nicht bewertbar\grqq{} und nicht als Null.

Ebenfalls nicht enthalten: Scope 2 und Scope 3. Dafür existiert bis heute keine freie Quelle mit brauchbarer Abdeckung. Die kalifornische Offenlegungspflicht ändert das ab November 2026 für Scope 1 und 2 und ab 2027 für Scope 3.

\section{Reproduktion}

\begin{center}\small
\begin{tabular}{lP{9.4cm}}
\toprule
\texttt{cp .env.example .env} & Schlüssel eintragen (zwei Stück, beide kostenlos) \\
\texttt{scripts/01\_fetch.py} & alle Quellen ziehen, als Parquet cachen \\
\texttt{scripts/02\_analyse.py} & Zuordnung, Monte Carlo, beide Perioden \\
\texttt{scripts/05\_neue\_dimensionen.py} & Abdeckung und Unabhängigkeit der Zusatzachsen \\
\texttt{scripts/07\_dataset.py} & gesammelten Datensatz und Web-Daten schreiben \\
\texttt{scripts/06\_katalog.py} & diesen Katalog erzeugen \\
\bottomrule
\end{tabular}
\end{center}

Einzelne Quelle neu ziehen: \texttt{scripts/01\_fetch.py <name> -{}-force}.

\end{document}
